\section{Classical inputs and the two-volume reduction}
\label{pre:section}

We work on
\[
 T=\mathbb R^3/(2\mathbb Z^3),
 \qquad V(E)=\frac{|E|}{8},
 \qquad A(E)=\frac{|\partial E|}{8}.
\]
Here and below perimeter is understood in the finite-perimeter sense; a
minimizer is represented by its regular region. The map from the unit torus
to \(T\) doubles lengths. Consequently \(V\) is the original volume fraction,
whereas \(A\) is one half of the original perimeter. The outward unit normal
is \(N\), and \(h=k_1+k_2\) is the sum of the principal curvatures in the
metric of \(T\), positive on a ball. It is not divided by eight.

Standard compactness and regularity for isoperimetric regions imply existence
at every \(0<V<1\), with smooth embedded boundary of constant mean curvature.
In a flat metric this boundary is real analytic: its local CMC graph
equation is analytic elliptic, so
\cite[Theorem~1.1]{Blatt} applies. The second variation is
nonnegative on all volume-preserving variations. We use the following
classical structural theorem in precisely its unrestricted, identity-group
case.

\begin{theorem}[Classical flat-manifold structure]\label{pre:structure}
Let \(E\) be an isoperimetric region in a complete orientable flat
three-manifold, with closed boundary \(\Sigma\).
Either \(\Sigma\) is a union of parallel planar surfaces, or \(\Sigma\)
is connected. In the connected case:
\begin{enumerate}
\renewcommand{\labelenumi}{\textup{(\roman{enumi})}}
\item genus zero gives a round sphere;
\item genus one gives a flat quotient of a plane or a circular cylinder;
\item the genus is at most four;
\item in genus two or three, \(H^2|\Sigma|\le 2\pi\), where \(H=h/2\);
in genus four, \(H=0\).
\end{enumerate}
\end{theorem}

These statements are Theorems~5 and~7 of Hauswirth--P\'erez--Romon--Ros
\cite{HPRR}, applied with the identity symmetry group. In particular the
orientation-preserving hypothesis of their Theorem~7 is satisfied. Their
genus-one assertion builds on the stability results of
Ritor\'e--Ros \cite{RitoreRos92,RitoreRos96}. Existence and regularity in the
compact torus setting are recorded in \cite[Theorem~3]{HPRR}.
We call a minimizing boundary \emph{standard} if it is planar, spherical,
or a circular-cylinder quotient. Every nonstandard minimizing boundary
is therefore connected of genus \(g\ge2\) and satisfies
\begin{equation}
 h^2 A\le\pi.
 \label{pre:curvature}
\end{equation}
Indeed \(h^2A=H^2|\Sigma|/2\), and the genus-four case has \(h=0\).

\subsection{The standard competitors}

For \(0\le V\le1\), set \(v=\min(V,1-V)\) and
\begin{equation}
 F(V)=\min\left\{
       (9\pi/2)^{1/3}v^{2/3},\ \sqrt{\pi v},\ 1
       \right\},
 \qquad V_1=\frac{4\pi}{81},\quad V_2=\frac1\pi.
 \label{pre:envelope}
\end{equation}
The first, second, and third expressions are the respective minima on
\(0\le v\le V_1\), \(V_1\le v\le V_2\), and \(V_2\le v\le1/2\).

\begin{lemma}\label{pre:standard}
The function \(F\) is attained by standard competitors. Every standard
minimizing boundary has \(A\ge F(V)\); equality gives exactly the balls,
shortest-geodesic tubes, and coordinate slabs in
Theorem~\ref{intro:main}, including the indicated ties and complements.
\end{lemma}

\begin{proof}
A sphere in \(T\) lifts to an embedded sphere in \(\mathbb R^3\).
Its distinct lattice translates have disjoint interiors: equal-radius
spheres with intersecting interiors would have intersecting boundaries.
Thus it bounds an embedded ball on one side. For its radius \(R\),
\[
 W=\frac{\pi R^3}{6},\qquad A=\frac{\pi R^2}{2}
       =(9\pi/2)^{1/3}W^{2/3},
\]
where \(W\) denotes the ball-side volume. Since \(W\ge v\), this area
is at least the first expression in \eqref{pre:envelope}.

For a compact circular-cylinder quotient, the stabilizer of its lifted
axis contains a nonzero axial lattice translation. Its primitive core
therefore has length \(2k\), with \(k=|m|\ge1\) for a primitive nonzero
\(m\in\mathbb Z^3\). Embeddedness gives a solid tube on one side, with
\[
 W=\frac{k\pi R^2}{4},\qquad A=\frac{k\pi R}{2}
       =\sqrt{\pi kW}\ge\sqrt{\pi v}.
\]
This comparison concerns the area of the given tube and does not require
a shortest tube of the possibly larger volume \(W\) to embed.
If such a boundary minimizes, \(A=I(V)\le F(V)\), where \(I\) is the
isoperimetric profile. Equality throughout forces \(k=1\), \(W=v\),
and \(v\) in the closed cylindrical branch. The corresponding sphere
comparison similarly forces \(W=v\) and the spherical branch.

A connected closed planar surface has a primitive integral normal
\(m\) and physical area \(4|m|\). Disjoint planar surfaces must be
parallel, and a single such torus does not separate \(T\). At least
two are required to bound a region. Their total normalized area is
therefore at least one, with equality only for two tori with
\(|m|=1\). Their intervening regions are exactly coordinate slabs.

Finally the branches in \eqref{pre:envelope} are realized by an embedded
ball, a shortest tube, and a coordinate slab, respectively. At the last
volume of its branch the ball radius is \(2/3\); the largest tube radius
in its branch is \(2/\pi\). Both are less than the transverse injectivity
radius one. The branch intersections are exactly \(V_1,V_2\).
Complements give the remaining volumes. This proves both attainment
and the asserted equality classification.
\end{proof}

\subsection{Upper supports for the profile}

\begin{lemma}\label{pre:supports}
The profile \(I:[0,1]\to[0,\infty)\) is continuous, concave, and satisfies
\(I(V)=I(1-V)\). If a minimizer at \(V\in(0,1)\) has sum curvature \(h\),
there is a smooth local upper support \(C\) for \(I\) at \(V\) such that
\(C'(V)=h\) and \(C''(V)\le0\). If its boundary is connected of genus \(g\),
then at contact
\begin{equation}
 C^2C''=-h^2C+\pi(1-g).
 \label{pre:parallel}
\end{equation}
Every such contact curvature is nonnegative for \(V<1/2\) and
nonpositive for \(V>1/2\).
\end{lemma}

\begin{proof}
Compactness of finite-perimeter sets gives lower semicontinuity of \(I\).
A local volume-changing smooth variation of a minimizer gives upper
semicontinuity at every interior volume. Since \(0\le I\le F\), continuity
also holds at the endpoints. Complementation gives the symmetry.

Let \(K=k_1k_2\) be the intrinsic Gauss curvature of \(\Sigma\).
Let \(a(\tau)\) and \(w(\tau)\) be normalized area and volume under
outward normal parallel displacement by \(\tau\). At \(\tau=0\),
\[
 w'=A,\qquad a'=hA,\qquad w''=hA,\qquad
 a''=\frac18\int_\Sigma 2K\,dS.
\]
Small parallels are embedded and \(w'>0\); expressing their area as a
function \(C(w)\) provides the upper support. The chain rule gives
\[
 C'=h,\qquad
 C^2C''=\frac18\int_\Sigma 2K\,dS-h^2A.
\]
For connected \(\Sigma\), Gauss--Bonnet gives \eqref{pre:parallel}.
For an arbitrary boundary,
\[
 \frac18\int_\Sigma(2K-h^2)\,dS
       =-\frac18\int_\Sigma(k_1^2+k_2^2)\,dS\le0,
\]
so \(C''\le0\). These are the parallel-variation formulas of
\cite[Section~3]{HPRR}, with the stated normalization.

For completeness, this contact property implies concavity without
differentiability of \(I\). Suppose \(I(c)<L(c)\) for a chord \(L\) on
\([a,b]\) and \(a<c<b\). Choose \(\varepsilon>0\) sufficiently small
that \(I(c)<L(c)+\varepsilon(c-a)(c-b)\).
The minimum of
\[
 I(x)-L(x)-\varepsilon(x-a)(x-b)
\]
is negative and occurs at an interior point \(x_0\).
Adding this minimum to the strictly convex quadratic produces a function
\(q\le I\) with equality at \(x_0\).
Any local upper support there satisfies \(C''(x_0)\ge q''=2\varepsilon\),
a contradiction.

For a concave function, the slope of a smooth upper support lies between
its one-sided derivatives:
\(I'_+(V)\le h\le I'_-(V)\).
The secant joining the equal values at \(V\) and \(1-V\) has slope zero.
Concavity therefore gives the asserted signs on either side of \(1/2\).
No curvature conclusion at exactly \(1/2\) is needed at this stage.
\end{proof}

\begin{proposition}[Two-volume reduction]\label{pre:reduction}
If every minimizer at \(V_1\) and \(V_2\) is standard, then
\(I=F\) on \([0,1]\), and every minimizer is one of the standard regions
of Lemma~\ref{pre:standard}.
\end{proposition}

\begin{proof}
The assumption, existence, and Lemma~\ref{pre:standard} give \(I=F\) at
\(V_1,V_2\), and hence also at \(1-V_2,1-V_1\).
On each open interval between consecutive points in
\[
 0,\ V_1,\ V_2,\ 1-V_2,\ 1-V_1,\ 1,
\]
the function \(F\) is smooth and
\begin{equation}
 F''+\frac{(F')^2}{F}\ge0.
 \label{pre:branch-ode}
\end{equation}
This follows directly for its power-law and constant branches.

If \(I-F\) had a negative minimum on one of the closed intervals,
it would occur in its interior. A minimizer there must be nonstandard
by Lemma~\ref{pre:standard}. For its upper support,
\(C=I<F\), \(C'=F'\), and \(C''\ge F''\) at contact. Hence
\[
 C''+\frac{(C')^2}{C}
       \ge F''+\frac{(F')^2}{F}\ge0,
\]
whereas \eqref{pre:parallel} makes the left-hand side
\(\pi(1-g)/C^2<0\). Thus \(I=F\) throughout.

Now consider any nonstandard minimizer at an interior point of one of
the smooth intervals. Its support touches \(F=I\) from above, so
\(C'=F'\) and \(C''\ge F''\), giving the same contradiction with \(C=F\).
The partition points are covered by the hypothesis and complementation.
In particular \(1/2\) belongs to the smooth constant interval, so it is
included. Lemma~\ref{pre:standard} completes the equality classification.
\end{proof}

The profile comparison and the classification inside each smooth branch
are Milman's Theorem~1.3 and Propositions~4.4--4.6 \cite{Milman}, based
on the differential method of Hauswirth--P\'erez--Romon--Ros.
The endpoint hypothesis here additionally excludes nonstandard
co-minimizers at the transitions. Establishing that hypothesis is the
task of the remaining geometric argument.
